% =====================================================================
%  Actuator Force Calculations — CIM + Koors40 + SFU1204 Ball Screw
%  CORRECTED: lead screw is SFU1204 (4 mm lead), not SFU1205 (5 mm).
%  The lead sits in the denominator of the core equation, so this
%  scales every downstream force by 5/4 = 1.25x.
%  Compile: pdflatex actuator-force-calcs.tex
% =====================================================================
\documentclass[11pt]{article}
\usepackage[margin=0.9in]{geometry}
\usepackage{amsmath,amssymb,booktabs,siunitx,array,bm}
\usepackage[table]{xcolor}
\setlength{\parindent}{0pt}
\renewcommand{\arraystretch}{1.15}

\begin{document}

\begin{center}
{\LARGE\bfseries Force Calculations --- CIM + Koors40 + SFU1204 Ball Screw}\\[4pt]
{\small Lead screw: \textbf{SFU1204, 4\,mm lead}}
\end{center}
\vspace{0.6em}

%% ---------------- given terms ----------------
\begin{center}
\begin{tabular}{@{}ll@{\hskip 2.2em}ll@{}}
\multicolumn{2}{@{}l}{\bfseries CIM Motor (am-0255)} & \multicolumn{2}{l}{\bfseries Koors40 Controller (am-5600)}\\
\midrule
Rated Voltage  & \SI{12}{\volt} DC          & Input Voltage Range   & \SI{6}{\volt} -- \SI{18}{\volt}\\
Free Speed     & \SI{5310}{rpm}             & Continuous Current    & \textbf{\SI{40}{\ampere}}\\
Free Current   & \SI{2.7}{\ampere}          & PWM Input (normal)    & \SI{1000}{\micro\second} -- \SI{2000}{\micro\second}\\
Stall Torque   & \textbf{\SI{2.425}{\newton\meter}} & PWM Input (extended) & \SI{500}{\micro\second} -- \SI{2500}{\micro\second}\\
Stall Current  & \SI{133}{\ampere}          & &\\
Peak Power     & \SI{337}{\watt}            & &\\
Shaft Diameter & \SI{8}{\milli\meter} (0.313 in.) & &\\
\bottomrule
\end{tabular}
\end{center}
\vspace{0.4em}

%% ---------------- core equation ----------------
\section*{Core Equation}

\begin{equation}
F_{\text{piston}} \;=\; T_{\text{motor}}\cdot GR \cdot \eta_{\text{gearbox}}\cdot
\frac{2\pi\,\eta_{\text{screw}}}{L}
\end{equation}

$\eta_{\text{gearbox}} = 0.85$,\quad $\eta_{\text{screw}} = 0.90$,\quad
$L = \SI{0.004}{\meter}$ \textbf{(SFU1204, 4\,mm lead)},\quad
$\eta_{\text{combined}} = 0.765$

\begin{equation}
F_{\text{piston}} \;=\; T_{\text{motor}}\cdot GR \times \mathbf{1201.7}
\qquad [\si{\newton} \text{ per } \si{\newton\meter}]
\end{equation}

\textbf{Operating Point A} --- Koors40 continuous limit (\SI{40}{\ampere}):
\[
T_{\text{motor}} = k_t\,(I - I_{\text{free}}) = 0.01823 \times (40 - 2.7) = \SI{0.680}{\newton\meter}
\]

%% ---------------- required vs available ----------------
\section*{Required Force at Master Cylinder}

Wilwood $\tfrac{5}{8}$\,in (\SI{15.875}{\milli\meter}) bore:
$A_{\text{piston}} = \tfrac{\pi}{4}(0.015875)^2 = \SI{1.979e-4}{\meter\squared}$

\begin{center}
\begin{tabular}{@{}llr@{}}
\toprule
\textbf{Pressure} & \textbf{Calc} & \textbf{Force}\\
\midrule
\SI{300}{psi}  (light)     & \SI{2.068}{\mega\pascal}  $\times A$ & \SI{409}{\newton}\\
\SI{600}{psi}  (moderate)  & \SI{4.137}{\mega\pascal}  $\times A$ & \SI{819}{\newton}\\
\SI{1000}{psi} (hard)      & \SI{6.895}{\mega\pascal}  $\times A$ & \SI{1365}{\newton}\\
\SI{1500}{psi} (near max)  & \SI{10.342}{\mega\pascal} $\times A$ & \SI{2047}{\newton}\\
\bottomrule
\end{tabular}
\end{center}

\section*{Available Force @ \SI{40}{\ampere}}

$F_{\text{piston}} = T_{\text{motor}}\cdot GR \times 1201.7$

\begin{center}
\begin{tabular}{@{}lrrrc@{}}
\toprule
\textbf{GR} & \textbf{$F_{\text{piston}}$} & \textbf{(lbf)} & \textbf{Pressure} & \textbf{$\geq$ 2047\,N}\\
\midrule
5:1              & \SI{4086}{\newton}  & 919  & \SI{2994}{psi}  & \checkmark\\
10:1             & \SI{8171}{\newton}  & 1837 & \SI{5987}{psi}  & \checkmark\\
\textbf{20:1}    & \textbf{\SI{16342}{\newton}} & \textbf{3674} & \textbf{\SI{11975}{psi}} & \checkmark\\
50:1             & \SI{40855}{\newton} & 9185 & \SI{29937}{psi} & \checkmark\\
\bottomrule
\end{tabular}
\end{center}

At the selected 20:1 ratio:
$\SI{16342}{\newton} \div \SI{1.979e-4}{\meter\squared} = \SI{11975}{psi}$,
i.e.\ \textbf{3674\,lbf} against a \textbf{1000\,lbf} operational target
(\textbf{3.7$\times$ margin}) and a \textbf{5000\,lbf} ball-screw axial rating.

%% ---------------- speed ----------------
\section*{Speed --- The Cost of the Reduction}

At 20:1 the screw turns at $5310/20 = \SI{265.5}{rpm}$:
\[
v = \frac{265.5}{60}\times\SI{4}{\milli\meter} = \SI{17.7}{\milli\meter\per\second}
\qquad\Rightarrow\qquad
\SI{25}{\milli\meter}\ \text{stroke} = \SI{1.41}{\second}
\]

A shorter lead trades speed for force: versus the 5\,mm assumption the actuator
is \textbf{25\% stronger and 20\% slower}. Acceptable, because a dynamometer
applies steady-state load during an RPM sweep rather than reacting to transients.

%% ---------------- what changed ----------------
\section*{Correction Summary}

\begin{center}
\begin{tabular}{@{}lrrr@{}}
\toprule
\textbf{Quantity} & \textbf{Was (5\,mm)} & \textbf{Now (4\,mm)} & \textbf{Change}\\
\midrule
Multiplier         & 961.3  & \textbf{1201.7} & $\times 1.25$\\
$F$ @ 5:1          & \SI{3268}{\newton}  & \SI{4086}{\newton}  & $\times 1.25$\\
$F$ @ 10:1         & \SI{6537}{\newton}  & \SI{8171}{\newton}  & $\times 1.25$\\
$F$ @ 20:1         & \SI{13074}{\newton} & \SI{16342}{\newton} & $\times 1.25$\\
$F$ @ 50:1         & \SI{32684}{\newton} & \SI{40855}{\newton} & $\times 1.25$\\
Pressure @ 20:1    & \SI{9580}{psi}      & \SI{11975}{psi}     & $\times 1.25$\\
Screw speed @ 20:1 & \SI{22.1}{\milli\meter\per\second} & \SI{17.7}{\milli\meter\per\second} & $\times 0.80$\\
\bottomrule
\end{tabular}
\end{center}

$L$ appears in the denominator, so a 20\% shorter lead raises every force by
exactly 25\% and lowers linear speed by exactly 20\%. No other term moves.

\end{document}
